% Phase-3 archived appendix unit
% Phase-2 Unit ID(s): D-03
% Original location: Relation knowledge / external evidence audit
% Original source SHA-256: 726916622e270ad0c5fd157c2a9821e84742065551d407a942ac397611d5460c
% Provenance: Post-hoc descriptive analysis
% Original label(s): none
% Compact PDF replacement: No grade taxonomy; factual evidence scope remains in Phase-3 Section E.2
% Content below is copied verbatim from the authoritative appendix.

% Original lines 830--839
The per-relation CAMELS ablation in
Table~\ref{tab:app-camels-per-relation} finds that the declared precipitation direction
accounts for essentially the joint content-sensitivity effect, whereas the PET-direction
contrast is small and not statistically separated from zero.
The audited scale mismatch — annual water-balance evidence applied as a daily constraint — is consistent with this interpretation but is not separately tested here.

Also, across the seven medical relations, the post-hoc literature-audit categories did not show a clear association with either content sensitivity or predictive utility.
Because the number of relations is small and the ordering of qualitative audit grades is not unique, we treat this comparison as descriptive rather than as evidence that stronger external support predicts larger downstream effects.
The CAMELS relations provide a separate qualitative example: the better-supported precipitation direction accounts for most of the observed content sensitivity, whereas the weaker PET direction contributes little, but this does not establish a general relationship between evidence strength and predictive effect.


